% =====================================================================
% Section IV -- Data and Semi-Synthetic Aggregates (2026-10-09, Alex: data and aggregate creation only, placed after Estimators).
% IV-A Source Datasets and IV-B CH and GB Datasets: draft 3 (reviewed 2026-10-02). IV-C: draft 2 (2026-10-02).
% The evaluation protocol and the pitfalls moved to their own section in main.tex. Sources of every number: paper/README.md.
% =====================================================================
\section{Data and Semi-Synthetic Aggregates}
\label{sec:data}

In this section, the source datasets and the two datasets built from them are described first. The construction of the semi-synthetic aggregates follows.

\subsection{Source Datasets}
\label{sec:datasets}

The source datasets used in this paper are listed in Table~\ref{tab:datasets}. Two of them provide submetered HP electricity, and two provide the load of households without HPs.

\begin{table*}[t]
  \centering
  \caption{Source datasets used in this paper. Counts are those reported by each dataset, before any selection.}
  \label{tab:datasets}
  \footnotesize
  \begin{tabular}{llllllp{3.3cm}}
    \toprule
    Dataset & Country & Households & HP measurement & Resolution & Period & Role \\
    \midrule
    HEAPO~\cite{brudermueller2025heapo} & Switzerland & \num{1408} & HP submeter (subset) & 15 min, daily & 2018--2024 & HP households, \dataCH \\
    Kaiser et al.~\cite{kaiser2026swiss} & Switzerland & \num{2447} & separate HP meter (subset) & 15 min & 2023--2024 & HP households and fillers, \dataCH \\
    EoH~\cite{esc2024eoh} & Great Britain & 742 & whole heating system & 30 min & 2020--2023 & HP households, \dataGB \\
    LCL~\cite{schofield2015lcl} & Great Britain & \num{4173}\textsuperscript{a} & none & 30 min & 2012--2014\textsuperscript{b} & fillers, \dataGB \\
    \bottomrule
    \multicolumn{7}{l}{\textsuperscript{a}Households on the flat-rate tariff in the files used here. \textsuperscript{b}Window used here.}\\
  \end{tabular}
  % TODO (07): add the transfer datasets (WPuQ, FeederBW, UKPN) once the transfer iteration fixes their use.
\end{table*}

The HEAPO dataset contains smart meter data of \num{1408} households with HPs in the canton of Zurich~\cite{brudermueller2025heapo}. It was recorded at 15-minute and daily resolution between 2018 and 2024, and a subset of the households has a separate submeter for the HP. Weather data from eight stations are included.

The dataset of Kaiser et al.\ contains 15-minute smart meter data of \num{2447} residential installations in Switzerland for 2023 and 2024~\cite{kaiser2026swiss}. Its metadata flag the major loads of each installation, such as HPs, backup heating rods, electric water heaters and storage heaters. Some HPs are measured by their own meter, next to the meter of the dwelling. The dataset has no weather data, so the MeteoSwiss station of Zurich-Kloten, the closest station to these installations, is used~\cite{kaiser2026swiss}.

The Electrification of Heat (EoH) Demonstration Project installed and monitored HPs in 742 homes in South East England, North East England and South East Scotland~\cite{esc2024eoh}. Its cleansed 30-minute dataset covers November 2020 to September 2023. The EoH data measure the electricity of the whole heating system, which includes the compressor, the backup heater, the immersion heater for hot water and the circulation pumps. The homes are grouped into weather groups, which share one outdoor temperature series. These groups are used as weather stations in this paper.

The Low Carbon London (LCL) trial recorded 30-minute smart meter data of London households~\cite{schofield2015lcl}. Part of these households were placed on a dynamic time-of-use tariff in 2013, while the others stayed on a flat-rate tariff. The LCL data contain no information on HPs. Outdoor temperature for these households is the daily mean at London Heathrow, obtained from Meteostat~\cite{meteostat}.

\subsection{CH and GB Datasets}
\label{sec:chgb}

Two datasets of HP households are built from these sources, the \dataCH{} and the \dataGB. Their main properties are summarised in Table~\ref{tab:chgb}. In both, the households without HPs are referred to as fillers, since they fill each aggregate up to its size.

\begin{table}[t]
  \centering
  \caption{The CH and GB datasets after selection. The HP channel is the measurement from which the HP labels are computed.}
  \label{tab:chgb}
  \footnotesize
  \begin{tabular}{>{\raggedright\arraybackslash}p{0.26\columnwidth}>{\raggedright\arraybackslash}p{0.27\columnwidth}>{\raggedright\arraybackslash}p{0.31\columnwidth}}
    \toprule
    & \dataCH & \dataGB \\
    \midrule
    HP households & 86 & 384 (319 in the replication year) \\
    Weather stations & 5 & 26 \\
    Period & Jan--Dec 2023 & Nov 2021--Oct 2022 \\
    Resolution & 15 min & 30 min \\
    HP channel & HP submeter & whole heating system \\
    Fillers & \num{1291} (Kaiser et al.) & \num{3199} (LCL) \\
    Filler period & same year & 2012--2014, matched by analog days \\
    \bottomrule
  \end{tabular}
\end{table}

\subsubsection{\dataCH}
Calendar year 2023 is used as the study period, since it is the only year covered by both Swiss datasets. A meter is kept only if at least 90\% of its raw 15-minute samples in that year are valid. This rule is checked before any gap is filled, so that interpolated values cannot inflate the coverage. Remaining gaps are then interpolated linearly in time.

Three types of HP household pass this rule. HEAPO contributes 47 households with an HP submeter and a channel for the remaining household load. The Kaiser et al.\ dataset contributes 24 dwellings whose meter is paired with a separate HP meter. It also contributes 15 HP meters in single-family houses that have no dwelling meter. Each of these 15 HP meters is given the load of one unused single-family dwelling without HP, drawn at random with a fixed seed. In total, the \dataCH{} has 86 HP households spread over five weather stations.

The fillers are the \num{1291} dwellings of the Kaiser et al.\ dataset whose metadata flag no Thermostatically Controlled Load (TCL). These dwellings have no HP, electric water heater, storage heater or direct electric heating. It should be noted that the own non-HP load of an HP household is kept as measured, and can include an electric water heater.

A smaller sensitivity subset keeps only the 47 HP households at the Zurich-Kloten station. % TODO: confirm whether this subset (repo name B) is reported in the paper (DECISIONS 2026-09-29: sensitivity arm).

\subsubsection{\dataGB}
A home of the EoH project is eligible if at least 90\% of its 30-minute bins are valid within a 12-month window. Gaps of up to two hours are interpolated, while longer gaps are filled with a donor day from the same home and day type. Runs of zero electricity of six hours or more are treated as missing only when the heat meter still records heat output, which indicates a meter dropout. Runs without heat output are kept as real switch-offs. Hybrid systems are excluded, since their boiler data are unreliable.

The 12-month window starts in the month that gives the most eligible homes. Among the start months from June 2021 to January 2022, this is November 2021, so the main window runs from November 2021 to October 2022. One home is excluded because its electricity reads close to zero on 15\% of the bins in which heat is delivered. This pattern points to a silent electricity meter rather than to real behaviour.

Weather groups with more than 5\% missing temperature are filled from the best-correlated group if that correlation is at least 0.98, and are dropped with their homes otherwise. The resulting \dataGB{} has 384 homes in 26 groups. Of these, 217 have an air-source HP, 153 a high-temperature air-source HP and 14 a ground-source HP.

A second window, from October 2022 to 28 September 2023, serves as a temporal replication. The end date is the last full day of EoH data. It contains 319 homes, of which 272 are also in the main window. It is therefore a second year of mostly the same homes, not an independent sample.

The EoH data contain no whole-house meter, so the fillers of the \dataGB{} come from the LCL trial. Only households on the flat-rate tariff are used, since the price signals of the dynamic tariff shift load in time. A household is kept if at least 90\% of its data are valid between July 2012 and February 2014. Households with any half-hour above 10~kWh, or with a run of zeros of 24 hours or more, are also removed. Of \num{4173} flat-rate households, \num{3199} remain.

The LCL data were recorded about nine years before the EoH data, so they are not used on their own calendar days. Instead, each filler is mapped to the EoH days by matching days with similar temperature and the same day type, as described in Section~\ref{sec:aggregates}.

\subsubsection{Labels per dataset}
The HP label of a household is computed from its HP channel, which differs between the two datasets. In the \dataCH{} the channel is the HP submeter only, whereas in the \dataGB{} it is the whole heating system, including the immersion heater. Labels are also computed at different resolutions, 15 minutes and 30 minutes respectively. For these reasons, results are compared between the two datasets through their gap to a physics-based estimator, and not through absolute errors.

% ---------------------------------------------------------------------
% III-C: draft 2 (2026-10-02, after Alex's review of draft 1: shorter, no diagnostic results, Paper A's notation).
% Sources of every number: paper/README.md. The detailed design envelope is in Appendix A.
% ---------------------------------------------------------------------
\subsection{Semi-Synthetic Aggregates}
\label{sec:aggregates}

Semi-synthetic aggregates, i.e.\ sums of measured households, are used as feeders with a known HP capacity. Their construction is shown in Fig.~\ref{fig:design}. Network losses are not modelled, since the aim is a realistic aggregate rather than an exact copy of a real feeder.

\begin{figure}[t]
\centering
\figplaceholder{Fig.~1 (to be drawn by Alex). Split of the households into training and test sets, then aggregates of $N$ consumers with $N_{hp}$ HPs of one weather station, built within each set.}
\caption{Construction of the semi-synthetic aggregates. Households are split into training and test sets before any aggregate is built.}
\label{fig:design}
\end{figure}

Households are split into a training set and a test set before any aggregate is built. One quarter of the HP households of each weather station is drawn for the test set, and the fillers are split in the same proportion. Each aggregate is then built from the households of one set only, so that no household appears in both.

Each aggregate has $N$ consumers, of which $N_{hp}$ have an HP. All its HPs come from one weather station, so that one temperature series drives its heating load. In the \dataCH, each HP household keeps its own non-HP load, and $N-N_{hp}$ fillers are added. In the \dataGB, the EoH homes have no record of their non-HP load, so each of the $N$ consumers receives one filler. All draws are made without replacement, so no household recurs within an aggregate.

The consumer count $N$ takes the values 10, 20, 40, 80 and 120. The HP penetration takes the values 5, 10, 20, 30, 50, 80 and 100\%. A combination of station, $N$ and penetration is kept only if its $N_{hp}$ HPs are at most three quarters of the HP households of that station in the set. Each kept combination holds ten training aggregates or five test aggregates. Per split, this gives 500 training and 135 test aggregates in the \dataCH. The \dataGB{} gives \num{2080} training and 455 test aggregates. The full design envelope is given in Appendix~\ref{app:envelope}.

\subsubsection{Fillers for the \dataGB{} from analog days}
\label{sec:analog}
The LCL fillers were recorded in 2012--2014 and the EoH homes in 2021--2023, so their calendar days cannot be aligned directly. Instead, every EoH day is paired with an LCL day of similar weather, called its analog day. The analog day has the same day type, and its day of the year lies within 30 days of that of the EoH day. Among these days, those within 1~K of the daily mean temperature of the EoH weather station are kept, and the window is widened up to 60 days if none is found. One of the three closest days in temperature is drawn at random, and all fillers of an aggregate take the same analog day.

This construction assumes that the load of a filler is independent of the HP load, given temperature, day type and time of year. This is the conditional independence assumption of statistical matching~\cite{dorazio2006statistical}. Choosing whole days by their nearest neighbours in temperature follows the resampling used in weather generators~\cite{lall1996nearest,rajagopalan1999knn} and the analog method of statistical downscaling~\cite{zorita1999analog}. Measured HP profiles have been added to national electricity demand~\cite{love2017addition}, and to household smart meter data~\cite{agrawal2026electricity}. To the best of the authors' knowledge, no previous work pairs HP and household loads from different years on temperature-analog days.

The temperature response of real households changes from one winter to the next. Rebuilding a held-out LCL winter from the other LCL days showed an error of the same size as this natural change. This variability limits any estimator at low HP penetration, where the fillers carry most of the temperature response. Penetrations up to 50\% are therefore marked as filler-variability-limited in every GB table. The same construction is also applied to the \dataCH, where the real aggregates are known, and its effect is reported next to the GB results.

The analog-day fillers cannot reproduce every property of a real feeder. The own non-HP load of an HP home is replaced by that of an unrelated household. The LCL loads also reflect the appliances and occupancy of 2012--2014. Finally, all flat-rate households are kept, so electrically heated LCL homes add to the temperature response of the fillers.
